% generated by roam.py handbook export -- не редактировать, править заметки
% 20260917205132-выпуклая_оболочка.org (body 23d1471820a4)

\begin{definition}[{Выпуклая оболочка}]\label[definition]{note:1AC657F9-2A3D-47D1-A71D-5B813D961260}
  Пусть $V$~--- векторное пространство, $M \subset V$,
  \[
  \mathrm{conv}(M) = \bigcap_Q \left\{ Q \subset V, Q \text{ -- выпуклое и } M \subset Q \right\}
  \] --- выпуклая оболочка
  Тут логика почти такая же, как с борелевской оболочкой:
  Выпуклая оболочка $M$~--- пересечение всех \hyperref[note:316B5BFA-8B9F-4267-9086-D6593AAB5F9B]{выпуклых множеств}, содержащих $M$.
\end{definition}
